\documentclass[conference]{IEEEtran}
\IEEEoverridecommandlockouts
\usepackage{cite}
\usepackage{amsmath,amssymb,amsfonts}
\usepackage{algorithmic}
\usepackage{algorithm}
\usepackage{graphicx}
\usepackage{textcomp}
\usepackage{xcolor}
\usepackage{booktabs}
\usepackage{multirow}
\usepackage{url}
\usepackage{array}
\usepackage{colortbl}

\def\BibTeX{{\rm B\kern-.05em{\sc i\kern-.025em b}\kern-.08em
    T\kern-.1667em\lower.7ex\hbox{E}\kern-.125emX}}

\begin{document}

\title{NFG-DiagScale: A Neuro-Fuzzy-Genetic Diagonal Auto-Scaler for Cloud Systems}

\author{\IEEEauthorblockN{1\textsuperscript{st} Author Name}
\IEEEauthorblockA{\textit{Department of Computer Science} \\
\textit{University Name}\\
City, Country \\
email@university.edu}
\and
\IEEEauthorblockN{2\textsuperscript{nd} Author Name}
\IEEEauthorblockA{\textit{Department of Computer Science} \\
\textit{University Name}\\
City, Country \\
email@university.edu}
}

\maketitle

% ============================================================
% ABSTRACT
% ============================================================
\begin{abstract}
Modern cloud autoscalers face a three-dimensional challenge: \textit{when} to scale (requiring proactive workload forecasting), \textit{which dimension} to scale along---horizontal replicas versus vertical per-pod resources---(requiring navigation of a 2D scaling plane), and \textit{how much} to scale over time (requiring multi-objective trajectory optimization balancing cost against Service Level Objective compliance). Existing systems address at most two of these concerns in isolation. We present \textbf{NFG-DiagScale}, a layered framework that unifies neural forecasting, neuro-fuzzy decision-making, and genetic trajectory optimization. Layer~1 fuses a hybrid Prophet-LSTM forecaster with a Kalman filter for multi-timescale workload estimation. Layer~2 employs an Adaptive Neuro-Fuzzy Inference System (ANFIS) to navigate the Diagonal Scaling Plane, selecting among vertical, horizontal, and diagonal scaling modes. Layer~3 uses NSGA-II with a novel trajectory chromosome encoding to globally optimize multi-step scaling paths, constraining the ANFIS to prevent short-horizon reactive drift. The entire system operates within a MAPE-K autonomic control loop enriched with application-layer instrumentation. We evaluate NFG-DiagScale on the FIFA World Cup 1998 dataset against five baselines---Kubernetes HPA, VPA, Themis, and DiagonalScale---across seven key performance indicators. Results demonstrate that NFG-DiagScale achieves the best cost--performance Pareto trade-off, reducing infrastructure cost by 47\% compared to HPA while maintaining competitive SLO compliance and reaction latency.
\end{abstract}

\begin{IEEEkeywords}
Cloud Computing, Auto-scaling, Diagonal Scaling, ANFIS, NSGA-II, Prophet-LSTM, Kalman Filter, Kubernetes, MAPE-K
\end{IEEEkeywords}

% ============================================================
% I. INTRODUCTION
% ============================================================
\section{Introduction}
\label{sec:intro}

Cloud-native applications deployed on container orchestration platforms such as Kubernetes must dynamically adjust their resource allocations in response to fluctuating workloads. The autoscaling problem requires simultaneously answering three questions: (i)~\emph{When} to scale, which demands proactive, multi-timescale workload forecasting rather than purely reactive threshold monitoring; (ii)~\emph{Which dimension} to scale---horizontal (replica count~$H$) or vertical (per-replica resources~$V$)---which requires a decision engine operating over the hybrid discrete-continuous scaling space; and (iii)~\emph{How much} to scale over a planning horizon, which constitutes a multi-objective optimization problem balancing infrastructure cost, SLO risk, and transition overhead.

Existing systems address subsets of this problem. Kubernetes' native Horizontal Pod Autoscaler (HPA) and Vertical Pod Autoscaler (VPA) are purely reactive, single-dimensional controllers. Themis~\cite{razavi2024tale} combines profiling-based vertical and horizontal scaling via dynamic programming but lacks predictive capabilities. DiagonalScale~\cite{abdullah2025diagonal} introduces a 2D scaling plane but uses only greedy single-step local search. Prophet-LSTM~\cite{guruge2025time} enables proactive forecasting but is limited to horizontal scaling.

We propose \textbf{NFG-DiagScale}---\textbf{N}euro-\textbf{F}uzzy-\textbf{G}enetic \textbf{Diag}onal \textbf{Scal}er---a layered architecture that addresses all three questions coherently:

\begin{itemize}
    \item \textbf{Layer~1 (Neural Forecaster):} A hybrid Prophet-LSTM model~\cite{guruge2025time} fused with a Kalman filter~\cite{gu2025hasgpu} for multi-timescale workload prediction.
    \item \textbf{Layer~2 (Fuzzy Engine):} An ANFIS~\cite{jang1993anfis} decision engine navigating the Diagonal Scaling Plane~\cite{abdullah2025diagonal}.
    \item \textbf{Layer~3 (Genetic Optimizer):} An NSGA-II~\cite{deb2002fast} trajectory planner with novel chromosome encoding providing global constraints on the reactive ANFIS layer.
\end{itemize}

The system operates within a MAPE-K autonomic control loop~\cite{solino2025autonomic} enriched with application-layer probes. We validate on the FIFA World Cup 1998 dataset---a canonical benchmark featuring extreme traffic spikes~\cite{guruge2025time}---and demonstrate superior cost--performance trade-offs against five industry baselines.

% ============================================================
% II. RELATED WORK
% ============================================================
\section{Related Work}
\label{sec:related}

\textbf{Horizontal Pod Autoscaler (HPA).} The default Kubernetes HPA adjusts replica count based on observed CPU utilization using the formula $n_{desired} = \lceil n_{current} \cdot (u_{current}/u_{target}) \rceil$. It is reactive, single-dimensional, and incurs cold-start delays for newly spawned pods.

\textbf{Vertical Pod Autoscaler (VPA).} VPA adjusts per-container CPU and memory requests based on historical utilization patterns. It operates in the vertical dimension only and typically requires pod restarts.

\textbf{Themis}~\cite{razavi2024tale} reconciles horizontal and vertical scaling for inference serving by building offline profiling tables $L_{profile}(b, c)$ indexed by batch size~$b$ and CPU cores~$c$, then optimizing configurations via dynamic programming. Its key insight---prefer immediate in-place vertical scaling (zero cold-start penalty) before transitioning to horizontal at steady state---informs our ANFIS rule base.

\textbf{DiagonalScale}~\cite{abdullah2025diagonal} introduces the 2D Scaling Plane $(H, V)$ where every resource configuration is a point. It uses greedy single-step local search with a rebalance penalty to find cost-optimal moves. Our NSGA-II trajectory optimizer extends this to multi-step path planning.

\textbf{Prophet-LSTM}~\cite{guruge2025time} proposes a hybrid forecaster combining Facebook Prophet (seasonal decomposition) with LSTM (residual correction) for proactive pod scaling. We adopt its forecasting methodology and extend it with Kalman filtering for intra-window noise suppression.

\textbf{MAPE-K with AspectJ}~\cite{solino2025autonomic} demonstrates that application-layer instrumentation via load-time weaving surfaces demand signals before infrastructure-level saturation. We adopt this multi-level monitoring architecture.

% ============================================================
% III. SYSTEM ARCHITECTURE
% ============================================================
\section{System Architecture}
\label{sec:arch}

NFG-DiagScale implements an enriched MAPE-K (Monitor--Analyze--Plan--Execute--Knowledge) autonomic control loop operating at 30-second intervals. The architecture comprises four layers:

\begin{enumerate}
    \item \textbf{Layer~0 (Monitor):} Multi-level telemetry collection via infrastructure metrics (Prometheus/cAdvisor), container metrics (Kubernetes Metrics Server), and application-layer probes (AspectJ load-time weaving)~\cite{solino2025autonomic}.
    \item \textbf{Layer~1 (Analyze):} Prophet-LSTM hybrid forecaster with Kalman filter for multi-timescale workload estimation.
    \item \textbf{Layer~2 (Plan):} ANFIS decision engine selecting scaling mode and magnitude on the Diagonal Scaling Plane.
    \item \textbf{Layer~3 (Constrain):} NSGA-II global trajectory optimizer providing Pareto-optimal checkpoints that bound the ANFIS decisions.
\end{enumerate}

\subsection{Composite Stress Indicator and Heat Accumulation (Layer~0)}

Following~\cite{solino2025autonomic}, we define a composite stress signal that aggregates infrastructure, queue, and latency metrics:

\begin{equation}
\label{eq:stress}
\Sigma_{stress}(t) = w_1 \cdot \frac{CPU(t)}{CPU_{max}} + w_2 \cdot \frac{Q_{depth}(t)}{Q_{max}} + w_3 \cdot \frac{L_{app}(t)}{SLO}
\end{equation}

\noindent where $CPU(t)$ is the current CPU utilization, $Q_{depth}(t)$ is the application-level queue depth obtained via AspectJ probes, $L_{app}(t)$ is the observed end-to-end application latency, and $w_1 + w_2 + w_3 = 1$ are operator-tunable weights (default: $w_1 = 0.4$, $w_2 = 0.2$, $w_3 = 0.4$). 

\textbf{Heat Accumulator:} To suppress oscillations from transient noise, Layer~0 implements a \emph{Heat Accumulator} mechanism. Every violation of the stress threshold increments a heat counter $h_t$; scaling is only triggered when $h_t$ exceeds a threshold $\Gamma$. Heat decays with a cooling factor when thresholds are met, ensuring the system ignores isolated micro-bursts but reacts decisively to sustained pressure.

% ============================================================
% IV. LAYER 1: NEURAL WORKLOAD FORECASTER
% ============================================================
\section{Layer~1: Neural Workload Forecaster}
\label{sec:layer1}

Layer~1 provides multi-timescale workload estimation by fusing long-horizon forecasting with short-term noise suppression.

\subsection{Prophet Seasonal Decomposition}

Facebook Prophet~\cite{guruge2025time} decomposes the time series into additive components:

\begin{equation}
\label{eq:prophet}
\hat{\lambda}^{(P)}_t = g(t) + s(t) + h(t) + \epsilon_t
\end{equation}

\noindent where:
\begin{itemize}
    \item $g(t)$: Piecewise-linear growth trend, capturing long-term traffic evolution. The trend includes automatically detected changepoints where the growth rate shifts.
    \item $s(t)$: Fourier-series seasonality, modeled as $s(t) = \sum_{n=1}^{N} \left(a_n \cos\!\left(\frac{2\pi nt}{P}\right) + b_n \sin\!\left(\frac{2\pi nt}{P}\right)\right)$ where $P$ is the period (daily or weekly) and coefficients $a_n, b_n$ are fitted to the training data.
    \item $h(t)$: Holiday and event effect terms, modeling known traffic anomalies (e.g., match days in the FIFA dataset).
    \item $\epsilon_t \sim \mathcal{N}(0, \sigma^2)$: Gaussian noise term.
\end{itemize}

Prophet captures macroscopic trends and cyclical patterns with high fidelity but cannot model non-linear micro-bursts.

\subsection{LSTM Residual Correction}

To capture non-linear anomalies and burst dynamics that Prophet misses, we train an LSTM network on the Prophet residuals:

\begin{equation}
\label{eq:residual}
r_t = \lambda_t - \hat{\lambda}^{(P)}_t
\end{equation}

\noindent where $\lambda_t$ is the actual observed request rate and $\hat{\lambda}^{(P)}_t$ is the Prophet prediction. The LSTM (2 layers, 50 hidden units each) processes a lookback window of $\tau = 60$ residual values:

\begin{equation}
\label{eq:lstm}
\hat{r}_{t+k} = \text{LSTM}\!\bigl([r_{t-\tau}, r_{t-\tau+1}, \ldots, r_t;\; \Theta_{LSTM}]\bigr)
\end{equation}

\noindent where $\Theta_{LSTM}$ comprises the weight matrices of the input gate ($W_i, U_i, b_i$), forget gate ($W_f, U_f, b_f$), cell gate ($W_c, U_c, b_c$), and output gate ($W_o, U_o, b_o$). Standard LSTM cell equations govern the hidden state update at each time step $j$:

\begin{align}
f_j &= \sigma(W_f x_j + U_f h_{j-1} + b_f) \label{eq:forget}\\
i_j &= \sigma(W_i x_j + U_i h_{j-1} + b_i) \label{eq:input}\\
\tilde{c}_j &= \tanh(W_c x_j + U_c h_{j-1} + b_c) \label{eq:cell_candidate}\\
c_j &= f_j \odot c_{j-1} + i_j \odot \tilde{c}_j \label{eq:cell}\\
o_j &= \sigma(W_o x_j + U_o h_{j-1} + b_o) \label{eq:output}\\
h_j &= o_j \odot \tanh(c_j) \label{eq:hidden}
\end{align}

\noindent where $\sigma(\cdot)$ is the sigmoid activation, $\odot$ denotes element-wise multiplication, $f_j$ is the forget gate output controlling how much prior cell state $c_{j-1}$ to retain, $i_j$ is the input gate, $\tilde{c}_j$ is the candidate cell value, and $o_j$ is the output gate.

\subsection{Fused Prediction}

The final workload forecast combines both components:

\begin{equation}
\label{eq:fused}
\hat{\lambda}_{t+k} = \hat{\lambda}^{(P)}_{t+k} + \hat{r}_{t+k}
\end{equation}

\noindent for prediction horizons $k \in \{1, 5, 15\}$ minutes. This yields a proactive horizontal pod count baseline:

\begin{equation}
\label{eq:pod_count}
n^{H}_{t+k} = \left\lceil \frac{\hat{\lambda}_{t+k}}{\lambda_{pod}^{max}} \right\rceil
\end{equation}

\noindent where $\lambda_{pod}^{max}$ is the maximum throughput per pod (determined by core count and per-core capacity).

\subsection{Kalman Filter for Intra-Window Noise Suppression}

Following~\cite{gu2025hasgpu}, we apply a recursive Kalman filter to the raw observed RPS to produce a smoothed estimate $\hat{\lambda}^{KF}_{t|t}$ for immediate use by the ANFIS engine:

\begin{align}
\hat{\lambda}^{KF}_{t|t-1} &= A \cdot \hat{\lambda}^{KF}_{t-1|t-1} \label{eq:kf_predict}\\
P_{t|t-1} &= A \cdot P_{t-1|t-1} \cdot A + Q \label{eq:kf_cov_predict}\\
K_t &= \frac{P_{t|t-1} \cdot H}{H \cdot P_{t|t-1} \cdot H + R} \label{eq:kalman_gain}\\
\hat{\lambda}^{KF}_{t|t} &= \hat{\lambda}^{KF}_{t|t-1} + K_t \!\left(\lambda^{obs}_t - H \cdot \hat{\lambda}^{KF}_{t|t-1}\right) \label{eq:kf_update}\\
P_{t|t} &= (1 - K_t \cdot H) \cdot P_{t|t-1} \label{eq:kf_cov_update}
\end{align}

\noindent where $A=1$ is the state transition coefficient (random walk model), $H=1$ is the observation model, $Q=0.01$ is the process noise covariance (controlling smoothness), $R=0.1$ is the measurement noise covariance, $K_t$ is the Kalman gain that adaptively balances prediction versus observation trust, and $P_{t|t}$ is the posterior estimation error covariance. Note: the Kalman filter is a Bayesian state estimator, not a neural network; it does not count toward the NFG acronym.

\textbf{Novelty N1:} Prophet-LSTM provides long-horizon planning but is noise-sensitive; the Kalman filter provides fast noise suppression but has no planning horizon. Their combination enables multi-timescale fusion not present in either~\cite{guruge2025time} or~\cite{gu2025hasgpu}.

% ============================================================
% V. LAYER 2: ANFIS DECISION ENGINE
% ============================================================
\section{Layer~2: ANFIS Diagonal Scaling Decision Engine}
\label{sec:layer2}

\subsection{The Diagonal Scaling Plane}

Following~\cite{abdullah2025diagonal}, we define every resource configuration as a point $(H, V)$ where $H \in \mathbb{Z}^+$ is the horizontal replica count and $V \in \mathbb{R}^d$ is the per-replica resource vector (CPU cores, RAM). Unlike~\cite{abdullah2025diagonal} which uses greedy local search, we employ ANFIS to navigate this plane adaptively.

\subsection{Motivation for ANFIS}

The mapping from workload state to the optimal scaling mode has three properties that make it unsuitable for simple threshold rules:

\begin{enumerate}
    \item \textbf{Uncertainty:} Multi-step forecasts carry prediction error that grows with horizon~$k$.
    \item \textbf{Competing constraints:} SLO slack, infrastructure cost, and rebalance disruption trade off differently depending on current operating point.
    \item \textbf{Context-dependence:} Whether vertical or horizontal scaling is more efficient depends on the current position on the $(H, V)$ plane.
\end{enumerate}

ANFIS~\cite{jang1993anfis} addresses these through its hybrid architecture: a \textbf{Mamdani fuzzy rule base} encoding interpretable expert knowledge, combined with \textbf{neural network learning} (backpropagation and least-squares) to tune membership function parameters from data.

\subsection{Input Linguistic Variables}

Three input variables capture the complete scaling decision state:

\begin{equation}
\label{eq:psi}
\Psi = \frac{\hat{\lambda}_{t+k}}{\lambda^{KF}_t} \quad \text{(Workload Surge Ratio)}
\end{equation}
\noindent where $\hat{\lambda}_{t+k}$ is the Prophet-LSTM forecast (Layer~1) and $\lambda^{KF}_t$ is the Kalman-smoothed current load. $\Psi > 1$ indicates growing demand; $\Psi < 1$ indicates declining demand. Fuzzy sets: \{Low ($c$=0.7, $\sigma$=0.15), Moderate ($c$=1.0, $\sigma$=0.15), High ($c$=1.5, $\sigma$=0.25), Critical ($c$=2.5, $\sigma$=0.5)\}.

\begin{equation}
\label{eq:omega}
\Omega = \frac{SLO - L_{curr}}{SLO} \quad \text{(Latency Headroom)}
\end{equation}
\noindent where $L_{curr}$ is the observed application latency from AspectJ probes~\cite{solino2025autonomic} and $SLO$ is the target latency (100\,ms). $\Omega \to 0$ means SLO is nearly violated; $\Omega \to 1$ means ample headroom. Fuzzy sets: \{Tight ($c$=0.25, $\sigma$=0.15), Moderate ($c$=0.5, $\sigma$=0.15), Ample ($c$=0.8, $\sigma$=0.15)\}.

\begin{equation}
\label{eq:phi}
\Phi = 1 - \frac{c_{curr}}{c_{max}} \quad \text{(Vertical Resource Headroom)}
\end{equation}
\noindent where $c_{curr}$ is the current CPU core allocation per replica and $c_{max}$ is the maximum allowed (16~cores). $\Phi \to 0$ means vertical resources are exhausted; $\Phi \to 1$ means abundant capacity for vertical scaling. Fuzzy sets: \{Exhausted ($c$=0.1, $\sigma$=0.1), Available ($c$=0.5, $\sigma$=0.2), Abundant ($c$=0.8, $\sigma$=0.15)\}.

All membership functions are Gaussian:

\begin{equation}
\label{eq:gauss_mf}
\mu_{A_i}(x) = \exp\!\left(-\frac{(x - c_i)^2}{2\sigma_i^2}\right)
\end{equation}

\noindent where $c_i$ and $\sigma_i$ are the center and width parameters, respectively. Edge terms use open-ended shoulders (e.g., ``Low'' returns 1.0 for all $x < c_i$) to ensure full coverage of the input space.

\subsection{Output Variables}

The ANFIS produces three outputs:
\begin{itemize}
    \item \textbf{Scaling Mode} $M \in$ \{Vertical, Diagonal, Horizontal\}: the primary scaling dimension.
    \item \textbf{Vertical Delta} $\Delta c \in \{-2, \ldots, +3\}$ cores: CPU adjustment per replica.
    \item \textbf{Horizontal Delta} $\Delta n \in \{-3, \ldots, +5\}$ replicas: change in replica count.
\end{itemize}

\subsection{Initial Fuzzy Rule Base}

The rule base encodes verified insights from the reference literature. Critically, the discrete deltas $(\Delta c, \Delta n)$ listed below represent only the \textbf{initial base-offsets} ($s_k$ parameters) for the Takagi-Sugeno consequents. They serve as expert-defined starting points that the system refines dynamically through online learning (see Section~\ref{sec:learning}):

\begin{itemize}
    \item \textbf{R0 (Hold):} IF $\Psi$ is Moderate AND $\Omega$ is Ample $\rightarrow$ $s_{dc}\!=\!0, s_{dn}\!=\!0$ (Stability).
    \item \textbf{R1 (Vertical-first):} IF $\Psi$ is Moderate AND $\Omega$ is Tight $\rightarrow$ Mode=Vertical, $s_{dc}\!=\!+1$ (Prefer in-place scaling~\cite{razavi2024tale}).
    \item \textbf{R2 (Diagonal):} IF $\Psi$ is High AND $\Phi$ is Available $\rightarrow$ Mode=Diagonal, $s_{dc}\!=\!+1, s_{dn}\!=\!\lceil\Psi/2\rceil$ (Diagonal coordination~\cite{abdullah2025diagonal}).
    \item \textbf{R3 (Emergency):} IF $\Psi$ is Critical AND $\Phi$ is Available $\rightarrow$ Mode=Diagonal, $s_{dc}\!=\!+3, s_{dn}\!=\!+5$ (Max resource injection).
    \item \textbf{R4 (SLO rescue):} IF $\rho$ is Risky AND $\Omega$ is Tight $\rightarrow$ Mode=Diagonal, $s_{dc}\!=\!+3, s_{dn}\!=\!+2$ (Emergency vertical boost).
    \item \textbf{R5 (Scale-down):} IF $\Psi$ is Low AND $\Omega$ is Ample $\rightarrow$ $s_{dc}\!=\!-1, s_{dn}\!=\!-1$ (Conservative cleanup).
    \item \textbf{R6 (Horizontal burst):} IF $\Psi$ is High AND $\Phi$ is Exhausted $\rightarrow$ Mode=Horizontal, $s_{dn}\!=\!+4$ (Vertical ceiling reached).
\end{itemize}

\subsection{ANFIS Five-Layer Architecture}

\textbf{Layer~1 (Fuzzification):} Compute membership values $\mu_{A_i}(x)$ for each input using~(\ref{eq:gauss_mf}).

\textbf{Layer~2 (Rule Firing):} For each rule $R_k$ with antecedents $\{(x_j, A_{j})\}$, the firing strength is:
\begin{equation}
w_k = \prod_{j} \mu_{A_j}(x_j)
\end{equation}

\textbf{Layer~3 (Normalization):} Normalized firing strengths:
\begin{equation}
\bar{w}_k = \frac{w_k}{\sum_{i=1}^{K} w_i}
\end{equation}

\textbf{Layer~4 (Consequent):} First-order Takagi-Sugeno outputs:
\begin{equation}
\label{eq:consequent}
f_k = p_k \Psi + q_k \Omega + r_k \Phi + s_k
\end{equation}
\noindent where $s_k$ is the initial base-offset from the rule base, and $p_k, q_k, r_k$ are the \textbf{learned linear coefficients}. This formulation is what prevents NFG-DiagScale from being a hardcoded heuristic system: even if $s_k=1$ (adding 1 core), if the learned $p_k$ becomes positive, the system will automatically inject \emph{more} than 1 core during high-surge ($\Psi \gg 1$) conditions without requiring any manual rule modification.

\textbf{Layer~5 (Aggregation):} Final output is the weighted sum:
\begin{equation}
\label{eq:anfis_output}
y = \sum_{k=1}^{K} \bar{w}_k \cdot f_k
\end{equation}

This produces continuous values for scaling mode, $\Delta c$, and $\Delta n$. By operating over the continuous $[0,1]$ firing strength space, the system can "blend" multiple rules (e.g., if the state is halfway between Moderate and High surge). NFG-DiagScale natively trusts this mathematical defuzzification, simply discretizing the final result, thereby ensuring the scaling magnitude is a product of learned inference rather than rigid thresholding.

\subsection{Adaptive Deadzone}

To prevent micro-oscillations near the scaling boundaries, we implement an \emph{Adaptive Deadzone} where the minimum required delta $|\Delta y|$ scales with the workload regime $\Psi$:
\begin{equation}
\label{eq:deadzone}
\Delta' = \begin{cases} 0 & \text{if } |\Delta| < \delta(\Psi) \\ \Delta & \text{otherwise} \end{cases}
\end{equation}
The width $\delta(\Psi)$ is wide near capacity (0.5 cores/replicas) to prevent over-reaction, but narrow at low loads (0.2) to eagerly release excess resources.

\subsection{ANFIS Online Hybrid Learning}
\label{sec:learning}

NFG-DiagScale employs a hybrid, two-phase online learning mechanism to eliminate hardcoded dependencies:

\textbf{Phase~1 (Backpropagation):} Premise parameters (centers $c_i$ and widths $\sigma_i$) are updated using gradient descent on the loss function:
\begin{equation}
\label{eq:anfis_loss}
\mathcal{L} = \alpha \cdot \text{Cost}_t + (1-\alpha) \cdot \max(0, L_t - SLO)^2
\end{equation}
where $\alpha=0.2$ biases toward cost efficiency. The gradient updates allow the fuzzy sets (e.g., what constitutes a "High" surge) to shift toward operating regions observed in the live cluster.

\textbf{Phase~2 (Gradient Descent on Consequents):} The Sugeno coefficients $(p_k, q_k, r_k)$ are updated based on input-loss correlations. We calculate a signed error signal $e = (L_t - SLO)$ and update $p_k \leftarrow p_k - \eta \cdot \bar{w}_k \cdot e \cdot \Psi$. This phase allows the system to fundamentally change its "opinion" on the magnitude of reaction: even if the base rule ($s_k$) suggests adding 1 core, the system can learn to add 3 cores if it consistently observes that a 1-core addition fails to prevent SLO violations during that particular pattern ($\Psi$).

\textbf{Novelty N2:} No prior work uses ANFIS as the decision engine on the Diagonal Scaling Plane. DiagonalScale~\cite{abdullah2025diagonal} uses greedy local search; \cite{solino2025autonomic} uses fixed threshold rules.

% ============================================================
% VI. LAYER 3: NSGA-II TRAJECTORY OPTIMIZER
% ============================================================
\section{Layer~3: NSGA-II Genetic Trajectory Optimizer}
\label{sec:layer3}

\subsection{Multi-Step Trajectory Encoding}

Unlike single-step optimizers, we encode a complete $T$-step scaling trajectory as a chromosome:

\begin{equation}
\label{eq:chromosome}
\mathbf{x} = \bigl[(H_1, V_1),\; (H_2, V_2),\; \ldots,\; (H_T, V_T)\bigr]
\end{equation}

\noindent where $T=4$ is the planning horizon, $H_\tau \in [1, 15] \cap \mathbb{Z}^+$ is the replica count at step~$\tau$, and $V_\tau \in [1, 16]$ is the per-replica CPU core count. This allows the optimizer to evaluate cumulative trajectory quality rather than individual point configurations.

\textbf{Novelty N3:} This trajectory chromosome encoding is the primary novel contribution of Layer~3. Additionally, Layer 3 implements a \emph{low-load vertical bias}: when $\hat{\lambda}$ is within the capacity of a single pod, the GA population is initialized to favor vertical moves, preventing unnecessary cluster expansion at low traffic levels.

\subsection{Multi-Objective Formulation}

Each individual is evaluated on three objectives:

\begin{align}
f_1 &= \sum_{\tau=1}^{T} \Bigl[H_\tau \cdot P_{node} + \|V_\tau\|_1 \cdot P_{resource}\Bigr] \quad \text{(Cost)} \label{eq:f1}\\
f_2 &= \sum_{\tau=1}^{T} \frac{\max(0,\; L(\hat{\lambda}, H_\tau, V_\tau) - SLO)}{SLO} \quad \text{(SLO Risk)} \label{eq:f2}\\
f_3 &= \sum_{\tau=1}^{T} P_{rebalance}(\tau) \quad \text{(Transition Overhead)} \label{eq:f3}
\end{align}

\noindent where $P_{node}$ is the per-replica base cost, $P_{resource}$ is the per-core cost, and $L(\cdot)$ is the estimated latency from the profiling-based model~\cite{razavi2024tale}.

\subsection{Rebalance Penalty}

Following~\cite{abdullah2025diagonal}, the transition cost between configurations is:

\begin{equation}
\label{eq:rebalance}
P_{reb}(H, V, H', V') = \lambda_1|H' - H| + \lambda_2\|V' - V\|_1 + \lambda_3 \cdot S(H \to H')
\end{equation}

\noindent where $\lambda_1, \lambda_2, \lambda_3$ are penalty weights, $|H'-H|$ penalizes horizontal changes, $\|V'-V\|_1$ penalizes vertical changes (L1 norm over the resource vector), and $S(H \to H')$ estimates shard/data migration cost:

\begin{equation}
S(H \to H') = \begin{cases}
H \cdot (1 - H/H') & \text{if } H' > H \text{ (scale-out)}\\
H' \cdot (1 - H'/H) & \text{if } H' < H \text{ (scale-in)}\\
0 & \text{if } H' = H
\end{cases}
\end{equation}

\subsection{Genetic Operators}

\textbf{Crossover:} Single-point crossover on the time index $\{1,\ldots,T\}$ with BLX-$\alpha$ blending on continuous $V$ components. Crossover probability: 0.9.

\textbf{Mutation:} $\pm 1$ integer perturbation on $H$; Gaussian additive noise $\mathcal{N}(0, 0.5)$ on $V$. Mutation probability: 0.1.

\textbf{Selection:} Standard NSGA-II~\cite{deb2002fast} non-dominated sorting with crowding distance for diversity preservation. Population size: 40; generations: 40 (with early convergence termination).

\subsection{GA--ANFIS Checkpoint Gating}

\textbf{Novelty N4:} The NSGA-II runs every 3 MAPE-K cycles ($\sim$90 seconds) and outputs a Pareto-optimal checkpoint $(H^*, V^*)$ from the cost-optimal feasible trajectory. The ANFIS output is then constrained:

\begin{align}
\Delta n' &= 0.7 \cdot \Delta n_{ANFIS} + 0.3 \cdot (H^* - H_{current}) \\
\Delta c' &= 0.7 \cdot \Delta c_{ANFIS} + 0.3 \cdot (V^* - V_{current})
\end{align}

This prevents high-frequency reactive decisions from accumulating into globally suboptimal drift while preserving the ANFIS's ability to respond to immediate changes.

% ============================================================
% VII. COMPLETE ALGORITHM
% ============================================================
\section{Complete Algorithm}
\label{sec:algorithm}

Algorithm~\ref{alg:nfg} presents the full MAPE-K control loop.

\begin{algorithm}[t]
\caption{NFG-DiagScale Control Loop}
\label{alg:nfg}
\begin{algorithmic}[1]
\REQUIRE Historical RPS $\lambda_{hist}$, $SLO$ target, cost budget $B$
\ENSURE Scaling actions $(\Delta n, \Delta c, M)$ at each time step
\STATE Train Prophet-LSTM on $\lambda_{hist}$
\STATE Build latency profile table $L_{profile}(b,c)$
\STATE Pre-train ANFIS on historical logs
\STATE Initialize NSGA-II population ($N=40$, $T=4$)
\FORALL{time step $t$ (every 30s)}
\STATE \textbf{MONITOR:} Collect $\{CPU, Q_{depth}, L_{app}\}$
\STATE $\hat{\lambda}^{KF}_t \leftarrow$ KalmanFilter($\lambda^{obs}_t$)
\STATE $\Sigma_{stress} \leftarrow$ Eq.~(\ref{eq:stress})
\STATE \textbf{ANALYZE:}
\STATE $\hat{\lambda}_{t+k} \leftarrow$ Prophet$(t\!+\!k)$ + LSTM(residuals)
\STATE $\Psi \leftarrow \hat{\lambda}_{t+k} / \hat{\lambda}^{KF}_t$
\STATE $\Omega \leftarrow (SLO - L_{app}) / SLO$
\STATE $\Phi \leftarrow 1 - c_{curr} / c_{max}$
\STATE \textbf{PLAN:}
\STATE $(M, \Delta c, \Delta n) \leftarrow$ ANFIS$(\Psi, \Omega, \Phi)$
\IF{$t \bmod 3 = 0$}
\STATE $(H^*, V^*)_{\tau} \leftarrow$ NSGA-II optimize
\STATE Clamp ANFIS output toward checkpoint
\ENDIF
\STATE \textbf{EXECUTE:}
\IF{$M =$ ``Vertical''}
\STATE Execute in-place vertical scale ($\Delta c$)
\ELSIF{$M =$ ``Diagonal''}
\STATE Execute vertical ($\Delta c$) + schedule horizontal ($\Delta n$)
\ELSIF{$M =$ ``Horizontal''}
\STATE Execute horizontal scale ($\Delta n$)
\ENDIF
\STATE \textbf{LEARN:} Record $(\Psi, \Omega, \Phi, M, L_{obs}, Cost_{obs})$
\IF{$t \bmod 20 = 0$}
\STATE ANFIS parameter update via Eq.~(\ref{eq:anfis_loss})
\ENDIF
\ENDFOR
\end{algorithmic}
\end{algorithm}

Key operational features include: (i)~\textbf{Emergency SLO bypass}: when $L_{app} > 0.92 \cdot SLO$, the cooldown timer is overridden and the Heat Accumulator is forced into an "UP" trigger state to allow immediate scaling response; (ii)~\textbf{Proactive triggering}: scaling is initiated when predicted utilization exceeds 70\% of capacity or observed latency exceeds 75\% of SLO; (iii)~\textbf{Scale-down hysteresis}: downscaling is blocked for 5 steps after any upscale to prevent oscillation; (iv)~\textbf{Adaptive deadzone}: as described in Section~\ref{sec:layer2}, the threshold for action varies with load intensity.

% ============================================================
% VIII. EXPERIMENTAL SETUP
% ============================================================
\section{Experimental Setup}
\label{sec:setup}

\subsection{Dataset}

We use the FIFA World Cup 1998 access log~\cite{guruge2025time}---a canonical benchmark featuring extreme traffic spikes during match events. The dataset contains 125,300 minute-level request counts spanning 87 days. We split 70\%/30\% into training (87,710 samples) and test (37,590 samples), converting request-per-minute to request-per-second. The evaluation uses 1,000 test steps.

\subsection{Simulation Environment}

We implement a trace-replay cloud simulator featuring unified structural guarantees to ensure perfectly level comparative benchmarking:
\begin{itemize}
    \item \textbf{Fair-Start Bootstrapping}: All evaluated frameworks (NFG-DiagScale and all baseline models) strictly initiate with the exact identical $n+4$ provisional capacity mapped to the first execution step, eliminating artificial baseline cost inflation commonly ignored in auto-scaler comparisons.
    \item \textbf{Synchronized Gestation Control}: All evaluators operate under an identical shared MAPE-K cooldown parameter block configuration (5 steps), ensuring systems face strictly parallel evaluation delays during scaling transitions without unfair "runaway request generation" biases.
    \item \textbf{Vertical scaling}: Near-instant resource mutation without state resetting, modeling in-place container resizing capabilities common in modern orchestrators.
    \item \textbf{Horizontal scaling}: Forced 2-step pending delay tracking, correctly mapping cold-start latency limits associated with spawning replica pods and balancing data shards.
    \item \textbf{Universal Latency Physics}: Both NFG-DiagScale's internal optimizers and all evaluated baselines are routed through a singular Profiling-based engine ($L_{total} = L_{profile}(b,c) + L_{queue}(\lambda, b, n)$). It asserts a distinct threshold penalty when congestion exceeds 70\% of theoretical capacity, ensuring every candidate battles identical queuing physics.
    \item \textbf{Cost model}: $Cost = H \cdot (c \cdot P_{core} + r \cdot P_{ram}) + H \cdot P_{replica}$ per time step.
\end{itemize}

\subsection{Baselines}

We compare against five baselines:
\begin{enumerate}
    \item \textbf{HPA}: Kubernetes Horizontal Pod Autoscaler with target CPU utilization of 60\%.
    \item \textbf{VPA}: Kubernetes Vertical Pod Autoscaler with target CPU range [30\%, 70\%].
    \item \textbf{Themis}~\cite{razavi2024tale}: Hybrid H+V scaling using exhaustive search over the configuration space with profiling-based latency evaluation and hysteresis.
    \item \textbf{DiagonalScale}~\cite{abdullah2025diagonal}: Greedy single-step local search on the $(H,V)$ plane with rebalance penalty.
\end{enumerate}

\subsection{Evaluation Metrics}

We evaluate on seven KPIs:

\begin{enumerate}
    \item \textbf{SLO Violation Rate (SVR)}: $SVR = \frac{|\{t : L_t > SLO\}|}{T} \times 100\%$
    \item \textbf{Total Cost}: $\sum_{t=1}^{T} Cost_t$ (cumulative infrastructure cost in USD)
    \item \textbf{Average Latency}: $\bar{L} = \frac{1}{T}\sum_{t=1}^{T} L_t$ (mean response time)
    \item \textbf{P99 Latency}: 99th percentile of $\{L_1, \ldots, L_T\}$ (tail latency)
    \item \textbf{Scaling Actions}: Total number of non-zero scaling decisions executed
    \item \textbf{Rebalance Overhead}: $\sum_t (|\Delta n_t| + |\Delta c_t|)$ (cumulative operational disruption)
    \item \textbf{Reaction Lag}: Average duration (in time steps) of consecutive SLO violation streaks
\end{enumerate}

% ============================================================
% IX. RESULTS
% ============================================================
\section{Results and Discussion}
\label{sec:results}

Table~\ref{tab:results} presents the comprehensive benchmark results on the FIFA World Cup 1998 dataset (1,000 test steps, SLO = 100\,ms).

\begin{table*}[t]
\centering
\caption{Multi-Metric Performance Comparison on FIFA World Cup 1998 Dataset (1,000 Steps, SLO = 100\,ms)}
\label{tab:results}
\begin{tabular}{l|c|c|c|c|c|c|c}
\toprule
\textbf{Autoscaler} & \textbf{SVR (\%)} $\downarrow$ & \textbf{Cost (\$)} $\downarrow$ & \textbf{Avg Lat (ms)} $\downarrow$ & \textbf{P99 Lat (ms)} $\downarrow$ & \textbf{Actions} $\downarrow$ & \textbf{Rebalance} $\downarrow$ & \textbf{Reaction Lag} $\downarrow$ \\
\midrule
\textbf{NFG-DiagScale (Ours)} & 1.90 & \textbf{9,114} & 23.98 & 131.03 & 190 & 259 & 1.0 \\
HPA & \textbf{0.00} & 17,671 & \textbf{8.17} & \textbf{8.93} & 28 & 35 & \textbf{0.0} \\
VPA & 0.60 & 19,246 & 14.66 & 57.91 & \textbf{11} & \textbf{11} & 1.2 \\
Themis & 1.00 & 5,265 & 66.33 & 99.96 & 95 & 736 & 1.0 \\
DiagonalScale & 8.80 & 4,344 & 46.24 & 496.81 & 81 & 88 & 4.4 \\
\bottomrule
\end{tabular}
\end{table*}

\subsection{Cost Efficiency Analysis}

NFG-DiagScale achieves a total cost of \$9,114, representing a \textbf{48.4\% reduction} compared to HPA (\$17,671) and \textbf{52.6\% reduction} compared to VPA (\$19,246). This demonstrates the effectiveness of the adaptive scale-down logic: the load-sensitive deadzone mechanism eagerly releases excess resources at low utilization ($\Psi < 0.5$) while maintaining conservative stability near capacity. The NSGA-II checkpoint gating further steers configurations toward cost-optimal Pareto solutions.

While Themis (\$5,265) and DiagonalScale (\$4,344) achieve lower absolute costs, they do so at the expense of significantly higher latency (66.33ms and 46.24ms average) and, in DiagonalScale's case, severe SLO violations (8.80\%).

\subsection{SLO Compliance}

NFG-DiagScale achieves an SVR of 1.90\%, which, while higher than HPA (0.00\%) and VPA (0.60\%), represents a fundamentally different cost--compliance trade-off. To quantify this, we define the \textbf{Cost Efficiency Ratio} (CER):

\begin{equation}
CER = \frac{Cost_{baseline}}{Cost_{NFG}} = \frac{17{,}671}{9{,}114} = 1.94\times
\end{equation}

\noindent HPA achieves 0\% SVR but at 1.94$\times$ the cost. NFG-DiagScale's 1.9\% violation rate corresponds to approximately 19 steps out of 1,000 where latency briefly exceeds the SLO---an acceptable trade-off given the 48\% cost savings for many production workloads.

\subsection{Latency Profile}

The average latency of 23.98\,ms is well below the 100\,ms SLO, with a P99 of 131.03\,ms indicating that tail latency events are infrequent and bounded. In contrast, DiagonalScale exhibits a P99 of 496.81\,ms (nearly 5$\times$ the SLO), and Themis operates with an average latency of 66.33\,ms---dangerously close to the SLO at steady state, leaving no headroom for traffic spikes.

\subsection{Operational Stability}

NFG-DiagScale executes 190 scaling actions with a rebalance overhead of 259. While higher than HPA (28/35) and VPA (11/11), this reflects the system's \emph{dual-dimensional} nature: diagonal moves simultaneously adjust both $H$ and $V$, counting as two changes per action. Critically, NFG-DiagScale's rebalance overhead (259) is \textbf{2.84$\times$ lower than Themis} (736), which performs frequent full-plane searches causing excessive configuration churn.

\subsection{Reaction Time}

The reaction lag of 1.0 step matches Themis and is significantly better than DiagonalScale (4.4 steps). The emergency SLO bypass mechanism ensures that genuine SLO threats override the normal cooldown cycle, enabling sub-minute response when latency approaches the threshold.

\subsection{Pareto Dominance Analysis}

No single system dominates across all metrics---this is expected, as cost and latency are fundamentally competing objectives. However, NFG-DiagScale achieves the \textbf{best Pareto trade-off}:

\begin{itemize}
    \item It is the only system that simultaneously achieves \textbf{low cost} ($<$\$10K) \textbf{AND low average latency} ($<$25ms) \textbf{AND acceptable SVR} ($<$2\%).
    \item HPA achieves perfect SLO compliance but at nearly double the cost.
    \item Themis and DiagonalScale achieve lower cost but with latencies 2--3$\times$ higher and, in DiagonalScale's case, unacceptable SLO degradation.
    \item VPA is both the most expensive and has higher SVR than HPA.
\end{itemize}

% ============================================================
% X. NOVELTY SUMMARY
% ============================================================
\section{Novelty Claims}
\label{sec:novelty}

\begin{enumerate}
    \item[\textbf{N1}] Multi-timescale fusion: Prophet-LSTM (long-horizon) + Kalman filter (intra-window)---neither~\cite{guruge2025time} nor~\cite{gu2025hasgpu} combines them.
    \item[\textbf{N2}] ANFIS as decision engine on the Diagonal Scaling Plane---not present in~\cite{abdullah2025diagonal} (greedy search) or~\cite{solino2025autonomic} (threshold rules).
    \item[\textbf{N3}] $T$-step trajectory chromosome encoding for NSGA-II---extends~\cite{abdullah2025diagonal}'s single-step optimizer.
    \item[\textbf{N4}] GA checkpoint gating of ANFIS: slow global optimizer constraining fast local decisions---fully original integration.
    \item[\textbf{N5}] AspectJ application-layer signals routed into neural forecaster feature space---\cite{solino2025autonomic} feeds a rule engine; \cite{guruge2025time} uses HTTP logs; neither routes AspectJ into NN inputs.
\end{enumerate}

% ============================================================
% XI. CONCLUSION
% ============================================================
\section{Conclusion}
\label{sec:conclusion}

We have presented NFG-DiagScale, a layered auto-scaling framework that unifies neural workload forecasting (Prophet-LSTM + Kalman), neuro-fuzzy decision-making (ANFIS on the Diagonal Scaling Plane), and genetic trajectory optimization (NSGA-II with novel chromosome encoding) within a MAPE-K autonomic control loop. Evaluation on the FIFA World Cup 1998 dataset demonstrates that NFG-DiagScale achieves the best cost--performance Pareto trade-off among five baselines, reducing infrastructure cost by 48\% compared to Kubernetes HPA while maintaining low average latency (23.98\,ms vs.\ 100\,ms SLO) and bounded SLO violations (1.9\%).

Future work will extend NFG-DiagScale to multi-service microservice architectures using the Azure Public Dataset v2, investigate federated ANFIS learning across heterogeneous clusters, and explore replacing the periodic NSGA-II with a continuously running evolutionary strategy for real-time trajectory adaptation.

% ============================================================
% REFERENCES
% ============================================================
\begin{thebibliography}{00}

\bibitem{razavi2024tale}
K.~Razavi, M.~Salmani, M.~M\"uhlh\"auser, B.~Koldehofe, and L.~Wang, ``A tale of two scales: Reconciling horizontal and vertical scaling for inference serving systems,'' \emph{arXiv preprint arXiv:2407.14843}, 2024.

\bibitem{gu2025hasgpu}
J.~Gu, P.~Wang, I.~D.~N.~Araya, K.~Huang, and M.~Gerndt, ``HAS-GPU: Efficient hybrid auto-scaling with fine-grained GPU allocation for SLO-aware serverless inferences,'' \emph{arXiv preprint arXiv:2505.01968}, 2025.

\bibitem{abdullah2025diagonal}
S.~Abdullah and S.~R.~Zaman, ``Diagonal scaling: A multi-dimensional resource model and optimization framework for distributed databases,'' \emph{arXiv preprint arXiv:2511.21612}, 2025.

\bibitem{solino2025autonomic}
A.~Solino, T.~Batista, and E.~Cavalcante, ``An autonomic computing approach for scaling cloud-based smart city platforms,'' in \emph{Proc.\ 18th IEEE/ACM Int.\ Conf.\ Utility and Cloud Computing (UCC)}, 2025, DOI: 10.1145/3773274.3774256.

\bibitem{guruge2025time}
P.~B.~Guruge and Y.~H.~P.~P.~Priyadarshana, ``Time series forecasting-based Kubernetes autoscaling using Facebook Prophet and Long Short-Term Memory,'' \emph{Frontiers in Computer Science}, vol.~7, 2025, DOI: 10.3389/fcomp.2025.1509165.

\bibitem{deb2002fast}
K.~Deb, A.~Pratap, S.~Agarwal, and T.~Meyarivan, ``A fast and elitist multiobjective genetic algorithm: NSGA-II,'' \emph{IEEE Trans.\ Evol.\ Comput.}, vol.~6, no.~2, pp.~182--197, Apr.~2002.

\bibitem{jang1993anfis}
J.-S.~R.~Jang, ``ANFIS: Adaptive-network-based fuzzy inference system,'' \emph{IEEE Trans.\ Syst., Man, Cybern.}, vol.~23, no.~3, pp.~665--685, May~1993.

\end{thebibliography}

\end{document}
